% I/O list: cabinet drawings vs the PLC mapping, channel by channel
\newcommand{\ok}{\textcolor{green!50!black}{OK}}
\newcommand{\nok}[1]{\textcolor{pidcheck}{\textbf{X} #1}}
\newcommand{\chk}[1]{\textcolor{orange!80!black}{\textbf{!} #1}}
\newcommand{\iolisthead}{%
  \toprule
  \textbf{Module / ch} & \textbf{Drawing: tag, description} & \textbf{Terminal} & \textbf{Cable (860-001)} &
  \textbf{PLC word.bit} & \textbf{PLC variable (GVLMain.)} & \textbf{Status} \\
  \midrule}
\newcommand{\iocols}{@{}p{24mm}p{66mm}p{17mm}p{38mm}p{16mm}p{78mm}p{95mm}@{}}

% =====================================================================================================  CommonRIO
\pidsheet{CAS-PID-900}{I/O list 1/2: CommonRIO -R1-001 -- drawing vs PLC}{I/O list, PLC/RIO1 cabinet (792.001 rev 1) vs FB.RIO.CommonRIO}{\pidrev}{%
  \node[anchor=north west, inner sep=0pt] at (0,242) {%
    \fontsize{6.6}{8}\selectfont
    \begin{tabular}{@{}p{24mm}p{66mm}p{17mm}p{38mm}p{16mm}p{78mm}p{95mm}@{}}
      \iolisthead
      -R1-003 652 & Orbit 864 x6 via hub -MB01 (OT01--06) & MB01 & -- (none in list) & in/out 0--11 & GVLIO.CommonRIO.Serial652 $\rightarrow$ Pen.Input.Orbit864Sensor[0..5] & \ok \\
      -R1-004 638 ch1 & FT61 counter, cylinder 1 & -XX05 1--3 & PLC-PEN01-FT61-KX01 4G & 12 / 13 & Outlet.Cylinder[0].Input.OilFlowStatusWord / OilFlowCounter & \ok \\
      -R1-004 638 ch2 & FT62 counter, cylinder 2 & -XX05 5--7 & PLC-PEN01-FT62-KX01 4G & 14 / 15 & Outlet.Cylinder[1].Input... & \ok \\
      -R1-005 638 ch1 & FT63 counter, cylinder 3 & -XX05 9--11 & PLC-PEN01-FT63-KX01 4G & 16 / 17 & Outlet.Cylinder[2].Input... & \ok \\
      -R1-005 638 ch2 & FT64 counter, cylinder 4 & -XX05 13--15 & PLC-PEN01-FT64-KX01 4G & 18 / 19 & Outlet.Cylinder[3].Input... & \ok \\
      -R1-007 455 ch1 & QT11 CO\textsubscript{2} inside pen & -XX01 1/2 & PLC-PEN01-QT11-KX01 & 20 & Pen.Input.CO2 & \ok{} (no range in PLC) \\
      -R1-007 455 ch2 & PT11 total gas pressure & -XX01 3/4 & PLC-PEN01-PT11-KX01 & 21 & Pen.Input.TotalGasPressure & \ok{} (no range in PLC) \\
      -R1-007 455 ch3/4 & spare & -XX01 5--8 & -- & 22 / 23 & -- & \ok \\
      -R1-008 455 ch1 & FT41 oxygen flow line 1 (symbol says FT61) & -XX01 9/10 & PLC-PEN01-FT61-KX01 3G & 24 & OxySystem[0].Input.OxygenFlow & \chk{drawing / cable label FT61} \\
      -R1-008 455 ch2 & FT42 oxygen flow line 2 & -XX01 11/12 & PLC-PEN01-FT42-KX01 & 25 & OxySystem[1].Input.OxygenFlow & \ok \\
      -R1-008 455 ch3 & FT43 oxygen flow line 3 & -XX01 13/14 & PLC-PEN01-FT43-KX01 & 26 & OxySystem[2].Input.OxygenFlow & \ok \\
      -R1-008 455 ch4 & PT40 oxygen supply pressure & -XX01 15/16 & PLC-PEN01-PT40-KX01 & 27 & Pen.Input.OxygenInjectionPressure & \chk{name: supply pressure, no range / alarm} \\
      -R1-009 455 ch1--4 & PT61 / PT62 / PT63 / PT64 cyl.\ 1--2 open / close & -XX01 17--24 & PLC-PEN01-PT61..64-KX01 & 28--31 & Outlet.Cylinder[0/1].Input.PressureOpen / Close & \ok{} (no range in PLC) \\
      -R1-010 455 ch1--4 & PT65 / PT66 / PT67 / PT68 cyl.\ 3--4 open / close & -XX01 25--32 & PLC-PEN01-PT65..68-KX01 & 32--35 & Outlet.Cylinder[2/3].Input.PressureOpen / Close & \ok{} (no range in PLC) \\
      -R1-011 455 ch1--4 & spare (cable list: PT69--PT612?) & -XX01 33--40 & PLC-PEN01-PT69..612-KX01 & 36--39 & -- & \chk{cables without drawing} \\
      -R1-012 455 ch1--4 & spare & -XX01 41--48 & -- & 40--43 & -- & \ok{} (A1 resolved rev B) \\
      -R1-017 1405 ch1 & spare & -XX03 1/2 & -- & 44.0 & Outlet.HydraulicPump.Input.StarterFault & \nok{used by the PLC, spare on drawing (A2)} \\
      -R1-017 1405 ch2--16 & spare & -XX03 3--32 & -- & 44.1--44.15 & -- & \ok \\
      -R1-014 563 ch1 & OV41 oxygen valve line 1 & -XX02 1/2 & PLC-PEN01-OV41-KX01 & out 20 & OxySystem[0].OxyValve.Output.RawPosition & \ok \\
      -R1-014 563 ch2 & OV42 oxygen valve line 2 & -XX02 3/4 & PLC-PEN01-OV42-KX01 & out 21 & OxySystem[1].OxyValve.Output.RawPosition & \ok \\
      -R1-015 563 ch1 & OV43 oxygen valve line 3 & -XX02 5/6 & PLC-PEN01-OV43-KX01 & out 22 & OxySystem[2].OxyValve.Output.RawPosition & \ok \\
      -R1-015 563 ch2 & spare & -XX02 7/8 & -- & out 23 & -- & \ok \\
      -R1-019 1515 ch1--8 & XV61..XV64 \_O / \_C via relays XK01--XK08 & -XX04 1--16 & PLC-PEN01-XV61..64-KX01 4G & out 24.0--24.7 & Outlet.Cylinder[0..3].Output.OpenValve / CloseValve & \ok \\
      -R1-020 1515 ch1 & spare (rev 0: XV65 shut-off valve) & -XX04 17/18 & PLC-PEN01-XV65-KX01 & out 24.8 & Outlet.Output.FailCloseValve ("utg\aa r") & \chk{cable listed, removed from drawing} \\
      -R1-020 1515 ch2 & spare & -XX04 19/20 & -- & out 24.9 & Outlet.HydraulicPump.Output.Start & \nok{not wired on drawing} \\
      -R1-020 1515 ch3 & spare & -XX04 21/22 & -- & out 24.10 & Outlet.HydraulicPump.Output.Reset & \nok{not wired on drawing} \\
      -R1-020 1515 ch4--8 & spare & -XX04 23--32 & -- & out 24.11--15 & -- & \ok \\
      \midrule
      \multicolumn{7}{@{}l}{\textbf{Not on the drawings, used or expected by the PLC / FD:} hydraulic tank fault LSL61/TSH61 (Outlet.Input.HydraulicOilTankFault, not mapped);
        emergency stop (GVLMain.EmergencyStop, OxySystem[i].Input.EmergencyStop, never written);} \\
      \multicolumn{7}{@{}l}{CommonRIO field power / fuse status (Status.FieldPowerOK / FieldFuseOK default TRUE); hydraulic pump starter fault and start / reset (above).} \\
      \bottomrule
    \end{tabular}};
  \node[anchor=north west, align=left, font=\sffamily\tiny, text width=150mm] at (0,52) {%
    \textbf{Process image check (750-362, complex and analog modules first, digital packed after).}
    Drawing rev 1: inputs 0--11 (652), 12--19 (2x 638), 20--43 (6x 455), 44 (1405) = 45 words; outputs 0--11, 12--19, 20--23 (2x 563),
    24 (2x 1515) = 25 words. PLC: read 45 words with DI at word 44, write 25 words -- matches (A1 resolved in rev B; until rev A the PLC
    read 41 words with the DI at word 40).};
}

% =====================================================================================================  InletRIO
\pidsheet{CAS-PID-901}{I/O list 2/2: InletRIO 1--3 -R1-100 -- drawing vs PLC}{I/O list, PLC/RIO1 and RIO2/3 cabinets vs FB.RIO.InletRIO}{\pidrev}{%
  \node[anchor=north west, inner sep=0pt] at (0,242) {%
    \fontsize{6.6}{8}\selectfont
    \begin{tabular}{@{}p{24mm}p{66mm}p{17mm}p{38mm}p{16mm}p{78mm}p{95mm}@{}}
      \iolisthead
      -R1-106 455 ch1 & LT0n level inside pen & -XX06 1/2 & RIOn-PEN01-LT0n-KX01 & 0 & Inlet[i].Input.LevelInsidePen & \ok \\
      -R1-106 455 ch2 & LT2n level outside pen & -XX06 3/4 & RIOn-PEN01-LT2n-KX01 & 1 & Inlet[i].Input.LevelOutsidePen & \ok \\
      -R1-106 455 ch3 & PT4n pressure after oxy pump (PT41/42/43) & -XX06 5/6 & RIOn-PEN01-PT4n-KX01 & 2 & OxySystem[i].Input.OxyPumpPressure & \ok \\
      -R1-106 455 ch4 & PT(n+3) pressure after injector (PT44/45/46) & -XX06 7/8 & RIOn-PEN01-PT4(n+3)-KX01 & 3 & OxySystem[i].Input.OxynatorOutletPressure & \ok{} (FD numbers differ) \\
      -R1-101 1405 ch1 & XF05 24 V field breaker tripped & internal & -- & 4.0 & Inlet[i].Input.FieldPower24VTripped & \ok \\
      -R1-101 1405 ch2 & XF05 load $>$ 80 \% & internal & -- & 4.1 & Inlet[i].Input.FieldPower24VHighLoad & \ok{} (no alarm in PLC) \\
      -R1-101 1405 ch3 & XF01 230 V breaker tripped & internal & -- & 4.2 & Inlet[i].Input.MainCircuitBreakerTripped & \ok{} (no alarm in PLC) \\
      -R1-101 1405 ch4 & UPS XT01 backup power on & internal & -- & 4.3 & Inlet[i].Input.BackupPowerOn & \chk{RIO2/3: wired to +24 V, always TRUE} \\
      -R1-101 1405 ch5 & UPS XT01 battery low / faulty & internal & -- & 4.4 & Inlet[i].Input.BatteryFault & \chk{RIO2/3: wired to +24 V, always TRUE} \\
      -R1-101 1405 ch6--16 & spare & -- & -- & 4.5--4.15 & -- & \ok \\
      -R1-104 1405 ch1 & LS7n1H water detected, chamber 1 & -XX07 1/2 & RIOn-PEN01-LS7n1H-KX01 & 5.0 & Inlet[i].Input.BilgePumpSwitchHigh[0] & \ok{} (no alarm in PLC) \\
      -R1-104 1405 ch2 & LS7n1HH high-high, chamber 1 & -XX07 3/4 & RIOn-PEN01-LS7n1HH-KX01 & 5.1 & Inlet[i].Input.BilgePumpSwitchHighHigh[0] & \chk{no alarm, no 2oo3 trip} \\
      -R1-104 1405 ch3 / ch4 & LS7n2H / LS7n2HH, chamber 2 & -XX07 5--8 & RIOn-PEN01-LS7n2H/HH-KX01 & 5.2 / 5.3 & ...SwitchHigh[1] / HighHigh[1] & \ok \\
      -R1-104 1405 ch5 / ch6 & LS7n3H / LS7n3HH, chamber 3 & -XX07 9--12 & RIOn-PEN01-LS7n3H/HH-KX01 & 5.4 / 5.5 & ...SwitchHigh[2] / HighHigh[2] & \ok \\
      -R1-104 1405 ch7 & QE7n bilge pump breakers & -XX07 13/14 & RIOn-PEN01-QE7n-KX01 & 5.6 & Inlet[i].Input.CircuitBreakerBilgePump & \ok \\
      -R1-104 1405 ch8 & PU41-QS oxy pump leakage (RIO1 only) & -XX07 15/16 & DC1-RIO1-KX01 & 5.7 & OxySystem[i].Input.OxyPumpLeakDetedtion & \chk{RIO2/3 spare; not used in logic} \\
      -R1-104 1405 ch9 & PU41-TS oxy pump temperature (RIO1 only) & -XX07 17/18 & DC1-RIO1-KX01 & 5.8 & OxySystem[i].Input.OxyPumpHighTemperature & \chk{RIO2/3 spare; not used in logic} \\
      -R1-104 1405 ch10--16 & spare & -XX07 19--32 & -- & 5.9--5.15 & -- & \ok \\
      -R1-102 1515 ch1 & reset XF05 & internal & -- & out 0.0 & Inlet[i].Output.ResetFieldPower24V & \ok{} (no logic writes it) \\
      -R1-102 1515 ch2--8 & spare & -- & -- & out 0.1--0.7 & -- & \chk{DC n expects relays XK1--3 for bilge contactors} \\
      \midrule
      ATV650 DC n & oxy pump drive, embedded Ethernet & Cat6 & NC-DCn-KY01 & 8501/8502, 3201--3211, 7121 & OxySystem[i].OxyPump (InletRIO[i] DriveMaster) & \chk{Ethernet port not drawn on 874.004} \\
      \bottomrule
    \end{tabular}};
  \node[anchor=north west, align=left, font=\sffamily\tiny, text width=150mm] at (0,95) {%
    \textbf{Process image check.} Drawing: inputs 1 x 455 (words 0--3) + 2 x 1405 (words 4, 5) = 6 words; outputs 1 x 1515 = 1 word (8 bits).
    PLC: FC23 read 26 words / write 21 words from address 0. \textcolor{pidcheck}{Check on the real coupler that reading and writing
    beyond its process image is accepted (the simulator serves 64 words).}\\[1mm]
    \textbf{Inlet n:} n = 1 InletRIO1 = -R1-100 in the PLC/RIO1 cabinet (792.001 sh.26--31); n = 2 RIO2 cabinet (792.002, one generic
    drawing); n = 3 RIO3 cabinet (no drawing; from the cable list). PLC index i = n $-$ 1.};
}
